\subsection{Fase 5: Integración del bot con Telegram}

El sistema se integra con Telegram como plataforma de comunicación entre el usuario y el asistente. Para ello, se utiliza la biblioteca python-telegram-bot, que permite configurar la aplicación, registrar los comandos y gestionar los mensajes recibidos. La autenticación se realiza mediante un token proporcionado por BotFather, el cual permite establecer la conexión con el bot.

Para gestionar las operaciones del sistema se utiliza asyncio, que facilita la ejecución de operaciones asíncronas. En particular, el comando /asistente utiliza asyncio.to_thread() para ejecutar la función del asistente en un hilo separado, evitando bloquear directamente el ciclo de eventos. Por otra parte, el manejador de mensajes de texto libre ejecuta la función responder() directamente.

Los comandos se registran mediante app.add_handler() y las clases CommandHandler y MessageHandler, que permiten asociar cada tipo de mensaje con su respectiva función. Finalmente, app.run_polling() inicia el proceso de recepción de mensajes y mantiene el bot en funcionamiento mientras se ejecuta el programa.

En la variable \textit{BIENVENIDA} se almacenan los mensajes que el usuario recibe por parte del bot a la hora de iniciar una conversacion, mientras que en la variable \textit{AYUDA} se almacenan los ejemplos de consultas que el usuario puede hacerle al bot con el objetivo de orientarlo sobre como describir la pelicula que desea encontrar. En \textit{metodo_usuario} se almacena el metodo de recomendacion que el usuario seleccione.

Posteriormente, se procede a crear los manejadores de comandos y mensajes, siendo esta la parte en la que se define la logica que debe seguir el bot ante las acciones del usuario, estas funciones se ejecutan cuando Telegram recibe un mensaje como por ejemplo \textit{/start} o \textit{/asistente}.

Este ultimo permite utilizar la generacion de respuestas mediante un modelo de lenguaje LLM, funcion que permite al usuario recibir una respuesta con lenguaje mas natural. Primero se obtiene la consulta del usuario y despues de verificar que esta no este vacia, se ejecuta el asistente mediante la funcion \textit{asyncio.to\_thread}, la cual realiza la recuperacion de las peliculas mediante embedding y utiliza los resultados obtenidos como contexto para que la LLM pueda generar una respuesta mas completa.

\begin{lstlisting}[language=Python, caption={~}]
ALGORITMO ASINCRONO cmd_asistente(update, context)

    // Obtener la consulta escrita por el usuario
    consulta <- unir(context.args, " ")
    ESPERAR update.message.reply_text("Pensando...")

    // Ejecutar el asistente en un hilo independiente
    // mediante asyncio.to_thread para no bloquear el bot
    respuesta <- ESPERAR asyncio.to_thread(
        asistente,
        consulta
    )

    // Enviar la respuesta generada al usuario por Telegram
    ESPERAR update.message.reply_text(respuesta)
FIN ALGORITMO
\end{lstlisting}

Tambien se le permite al usuario poder procesar mensajes de texto libre sin necesidad de estar limitado a escoger alguno de los comandos con la funcion \textit{async def on_mensaje}, la cual se encarga de registrar el mensaje del usuario y procesarlo mediante el metodo de recomendacion que este haya seleccionado, devolviendo la ficha de la pelicula recomendada.

Se construye la aplicacion de Telegram mediante \textit{app = ApplicationBuilder().token(TOKEN).build()}, en donde se registran los manejadores. Se utilizan tres elementos clave, \textit{ApplicationBuilder()}, que inicia la configuracion de la aplicacion, \textit{token(TOKEN)}, que establece el token de autenticacion que fue proporcionado por Telegram con BotFather y \textit{build()}, que construye la aplicacion con la configuracion ya establecida.

Los comandos se registran con \textit{app.add_handler()}, cada \textit{CommandHandler} relaciona un comando con su respectiva funcion, mientras que los registros de los mensajes de texto libre se hacen con \textit{MessageHandler}.

Por ultimo, el bot se ejecuta con \textit{app.run\_polling()}, el cual se encarga de poner en funcionamiento el chatbot.